\documentclass[12pt]{article}

\usepackage{amssymb,amsmath,amsthm}

% gaya teorema (saya ingin semuanya menggunakan pencacah yang sama)
\theoremstyle{plain}
\newtheorem{theorem}{Teorema}[section]
\newtheorem{lemma}[theorem]{Lema}
\newtheorem{proposition}[theorem]{Proposisi}
\newtheorem{corollary}[theorem]{Korolari}
\newtheorem{conjecture}[theorem]{Konjektur}
\newtheorem{exercise}[theorem]{Latihan}
\theoremstyle{definition}
\newtheorem{definition}[theorem]{Definisi}
\newtheorem{example}[theorem]{Contoh}
\theoremstyle{remark}
\newtheorem{remark}[theorem]{Catatan}

% beberapa pengaturan penomoran
\numberwithin{equation}{section}

\usepackage{fancyhdr}
\usepackage{lastpage}
\setlength{\headheight}{15.2pt}
\renewcommand{\footrulewidth}{0.4pt}% nilai bawaan adalah 0pt
\setlength{\footskip}{30pt}
\pagestyle{fancy}

\makeatletter
\let\ps@plain\ps@fancy 
\makeatother

\lhead{MATH 742}
\chead{Representasi Terinduksi}
\rhead{Musim Semi 2023}
\lfoot{Terakhir Direvisi: \today}
\cfoot{}
\rfoot{\thepage\ dari \pageref{LastPage} }

\begin{document}

%%%%%%%%%%%%%%%%%%
%%%%%%%%%%%%%%%%%%
\title{Representasi Terinduksi}
\author{Alexander Duncan}

\maketitle

``Dasar-dasar'' representasi terinduksi dapat ditemukan dalam hampir
setiap buku pengantar teori karakter. Misalnya:
\cite[\S{16}]{AlperinBell},
\cite[\S{5.8-5.11}]{Etingof},
\cite[\S{3.3}]{FultonHarris},
\cite[\S{XVIII.6--9}]{Lang}, atau
\cite[\S{3.3}]{Serre}.
Namun, di sini kita membahas sedikit lebih jauh daripada ``dasar-dasarnya'';
lihat \cite[\S{7--10}]{Serre}.

%%%%%%%%%%%%%%%%%%
%%%%%%%%%%%%%%%%%%

\section{Aksi Grup}

Misalkan $G$ adalah grup.

\begin{definition}
Suatu \emph{himpunan-$G$} (kiri) adalah himpunan $X$ yang dilengkapi aksi (kiri) dari
grup hingga $G$. Untuk himpunan-$G$ $X$ dan $Y$, suatu \emph{morfisme himpunan-$G$}
atau \emph{pemetaan $G$-ekuivarian} adalah fungsi $f : X \to Y$
sedemikian sehingga $f(gx)=gf(x)$ untuk semua $x \in X$ dan $g \in G$.
Suatu \emph{isomorfisme himpunan-$G$} adalah bijeksi $G$-ekuivarian.
\end{definition}

Misalkan $H$ adalah subgrup dari $G$. Ingat bahwa suatu \emph{koset kiri dari $H$ dalam
$G$} adalah subhimpunan $G$ berbentuk
\[
gH = \{ gh \mid h \in H \}
\]
untuk suatu $g \in G$. Ingat bahwa koset-koset kiri mempartisi $G$.
Dengan $G/H$, kita menotasikan himpunan koset kiri dari $H$ dalam $G$.
Kita \emph{tidak} mengharapkan $G/H$ memiliki struktur grup, kecuali jika $H$
normal.

\begin{proposition}
Himpunan $G/H$ memiliki aksi kiri transitif kanonik dari $G$.
\end{proposition}

\begin{proof}
Aksi tersebut didefinisikan oleh
\[
g \cdot (aH) := (ga)H
\]
untuk $g,a \in G$.
Misalkan $aH=bH$ untuk $a,b \in G$. Maka $a=bh$ untuk suatu $h \in H$.
Jadi, $ga=gbh$ untuk semua $g \in G$. Dengan demikian, $gaH=gbH$.
Kita menyimpulkan bahwa aksi tersebut terdefinisi dengan baik.
Pemeriksaan dua aksioma aksi grup yang tersisa diserahkan sebagai latihan.
\end{proof}

\begin{remark}
Definisi aksi kiri pada $G/H$ lebih halus daripada yang mungkin tersirat oleh
pembacaan sepintas atas pembuktian tersebut.
Kita \emph{tidak} memiliki aksi kanan yang analog pada $G/H$
karena aksi tersebut tidak terdefinisi dengan baik. Syarat agar aksi kanan itu
terdefinisi dengan baik ekuivalen dengan syarat bahwa $H$ merupakan subgrup normal.
\end{remark}

\begin{exercise}
Misalkan $H,K$ adalah subgrup dari $G$.
Himpunan-$G$ $G/H$ dan $G/K$ isomorfik jika dan hanya jika
$H$ dan $K$ merupakan subgrup-subgrup yang berkonjugasi dalam $G$.
\end{exercise}

Ingat bahwa himpunan-$G$ $X$ disebut \emph{transitif} jika $X$ tak kosong
dan, untuk semua $x \in X$ dan
$y \in Y$, terdapat $g \in G$ sedemikian sehingga $gx=y$.

\begin{exercise}
Setiap himpunan-$G$ dapat ditulis secara tunggal sebagai gabungan saling lepas dari
himpunan-$G$ transitif.
\end{exercise}

\begin{proposition}
Setiap himpunan-$G$ transitif tak kosong $X$ isomorfik dengan $G/H$
untuk suatu subgrup $H \subseteq G$.
\end{proposition}

\begin{proof}
Karena $X$ tak kosong, kita dapat memilih suatu $x \in X$.
Misalkan $H$ adalah penstabil $x$ dalam $G$:
\[
H = \operatorname{Stab}_G(x) = \{ g \in G \mid gx=x \}.
\]
Kita mengklaim bahwa terdapat fungsi $G$-ekuivarian yang tunggal
$f : X \to G/H$ sedemikian sehingga $f(x)=H$.
Memang, karena $G$ bertindak secara transitif pada $X$,
syarat $f(gx)=gf(x)$ menentukan nilai $f$ di setiap titik
$X$. Untuk melihat bahwa fungsi ini terdefinisi dengan baik, cukup amati bahwa
$f(gx)=f(g'x)$ ekuivalen dengan $g^{-1}g'x=x$, yang ekuivalen dengan
$g^{-1}g' \in H$.
\end{proof}

\subsection{Aksi Kanan}

Semua pernyataan di atas tentang aksi grup \emph{kiri} memiliki
perumuman alami untuk aksi grup \emph{kanan}, yang sebagian besar kita serahkan sebagai
latihan bagi pembaca.
Kita menggunakan notasi $H \backslash G$ untuk himpunan koset kanan dari $H$
dalam $G$.

Latihan-latihan berikut sangat disarankan jika Anda belum pernah
memikirkan perbedaan antara aksi kiri dan kanan:

\begin{exercise}
Jika $X$ merupakan himpunan-$G$ \emph{kiri} dan $G$ abelian, terdapat
aksi \emph{kanan} pada $X$ melalui $x \cdot g := gx$ untuk $g \in G$ dan $x \in X$.
\end{exercise}

\begin{exercise}
Jika $X$ merupakan himpunan-$G$ \emph{kiri}, kita memperoleh aksi \emph{kanan}
pada $X$ melalui $x \cdot g := g^{-1}x$ untuk $g \in G$ dan $x \in X$.
\end{exercise}

\begin{exercise}
Fungsi $f : G/H \to H\backslash G$ yang diberikan oleh $f(gH)=Hg$ merupakan
bijeksi yang terdefinisi dengan baik jika dan hanya jika $H$ normal dalam $G$.
\end{exercise}

\subsection{Koset Ganda}

Misalkan $H, K$ adalah subgrup dari $G$.

\begin{definition}
Suatu \emph{koset ganda (H,K) dari G} adalah himpunan berbentuk
\[
HgK := \{ hgk \mid h \in H, k \in K \}
\]
untuk suatu $g \in G$.
Himpunan koset ganda dari $G$ dinotasikan dengan
$H \backslash G / K$.
\end{definition}

Ada beberapa sudut pandang ekuivalen mengenai koset ganda.
Koset ganda adalah kelas-kelas ekuivalensi dari relasi ekuivalensi
dengan $g \sim hgk$ untuk semua $h \in H$ dan $k \in K$.
Koset ganda adalah orbit-orbit aksi kiri $H$ pada $G/K$.
Koset ganda adalah orbit-orbit aksi kanan $K$ pada $H\backslash G$.
Terakhir, koset ganda adalah orbit-orbit aksi \emph{kiri}
dari $H \times K$ pada $G$ melalui $(h,k)\cdot g := hgk^{-1}$.

Masing-masing sudut pandang ini berguna dalam konteks yang berbeda.
Salah satu konsekuensi langsungnya ialah koset-koset ganda mempartisi $G$;
dengan kata lain, $HgK$ dan $Hg'K$ saling lepas atau sama,
dan setiap unsur $G$ termuat dalam suatu koset ganda.
Konsekuensi lain ialah setiap koset ganda sekaligus merupakan gabungan dari
koset-koset kanan $H$ dalam $G$ dan gabungan dari koset-koset kiri $K$ dalam $G$.

Perhatikan bahwa, tidak seperti koset biasa, semua koset ganda tidak selalu
memiliki kardinalitas yang sama. Meskipun demikian, kardinalitasnya mudah dihitung
dengan teorema orbit--penstabil:

\begin{exercise}
Jika $G$ hingga, maka
\[
|HgK| = \frac{|H| |K|}{|H \cap g K g^{-1}|}=
\frac{|H| |K|}{|g^{-1}Hg \cap K|}.
\]
\end{exercise}

\section{Definisi Representasi Terinduksi}

Misalkan $H$ adalah subgrup dari grup hingga $G$.

Untuk suatu representasi berdimensi hingga
$(V,\rho)$ dari $G$, kita mendefinisikan \emph{restriksi ke
$H$} sebagai representasi $(V,\rho|_H)$, dengan ruang vektor dasar
yang sama, tetapi homomorfisme grup $\rho|_H : H \to
\operatorname{GL}(V)$ memiliki domain $H$.
Kita sering menulis $\operatorname{Res}_H^G(V)$ atau $V {\downarrow}^G_H$
untuk menekankan aksi $H$, meskipun ruang vektornya tetap sama.

Representasi terinduksi memungkinkan kita bergerak ke arah sebaliknya.
Kita mulai dengan representasi dari $H$ dan menghasilkan representasi dari
$G$.
Representasi terinduksi dapat didefinisikan dengan mudah melalui sifat
universalnya:

\begin{definition}
Misalkan $H$ adalah subgrup dari grup hingga $G$ dan
$W$ suatu representasi berdimensi hingga dari $H$.
\emph{Representasi terinduksi} adalah representasi
$\operatorname{Ind}_H^G W$ dari $G$, bersama suatu
pemetaan $H$-ekuivarian $\iota : W \to \operatorname{Ind}_H^G W$ sedemikian sehingga,
untuk setiap representasi lain $V$ dari $G$ dan pemetaan $H$-ekuivarian
$\varphi : W \to V$, terdapat pemetaan $G$-ekuivarian tunggal
$\psi: \operatorname{Ind}_H^G W \to V$ sedemikian sehingga
$\psi \circ \iota = \varphi$.
\end{definition}

Notasi alternatif untuk $\operatorname{Ind}_H^G W$ adalah $W
{\uparrow}_H^G$.

Seperti biasa, sifat universal ringkas untuk dinyatakan dan berguna
secara teoretis, tetapi keberadaannya tidak langsung terlihat.
Di sisi lain, suatu konstruksi eksplisit dapat memuat berbagai
``detail implementasi'' yang tidak penting, berbeda-beda dari satu penulis ke penulis lain,
dan tidak bersifat mendasar.
Walaupun demikian, konstruksi langsung tetap harus diberikan dan kita
akan memerlukannya nanti.

\begin{proposition} \label{prop:ind_construction}
Representasi terinduksi ada dan tunggal hingga isomorfisme tunggal.
\end{proposition}

\begin{proof}
Ketunggalan mengikuti dari sifat pemetaan universal melalui
argumen kategori abstrak. Bagian yang menarik adalah keberadaan.

Misalkan diberikan grup hingga $G$ dengan subgrup $H$ dan representasi
$(W,\sigma)$ dari $H$.
Pilih suatu sistem wakil berbeda $X$ untuk $G/H$.
Misalkan $V$ adalah jumlah langsung $\bigoplus_{x \in X} W$.
Dengan melabeli suku langsung ke-$x$ dari $V$ sebagai $xW$,
kita memperoleh isomorfisme $\phi_x : W \to xW$.

Sekarang perhatikan bahwa
\[
\operatorname{End}_k(V) = \bigoplus_{x \in X} \bigoplus_{y \in X}
\operatorname{Hom}_k(xW,yW)
\]
memiliki struktur ``matriks blok''.
Untuk suatu unsur $\psi \in \operatorname{End}_k(V)$,
notasikan dengan $\psi^{y,x}$ restriksinya ke setiap
$\operatorname{Hom}_k(xW,yW)$.

Tetapkan $g \in G$; kita akan mendefinisikan suatu pemetaan
$\rho_g \in \operatorname{End}_k(V)$ dengan mendefinisikan setiap komponen
$\rho_g^{y,x}$.
Untuk semua $g \in G$ dan $x,y \in X$, kita definisikan
\[
\rho_g^{y,x} =
\begin{cases}
\phi_y \circ \sigma_{y^{-1}gx} \circ \phi_x^{-1} & \textrm{jika $gxH=yH$,}\\
0 & \textrm{selain itu},
\end{cases}
\]
dan perhatikan bahwa $y^{-1}gx \in H$, sehingga evaluasi $\sigma$
bermakna.

Memeriksa bahwa $(V,\rho)$ adalah representasi linear dari $G$
yang memenuhi sifat universal yang diinginkan memang membosankan, tetapi langsung.
\end{proof}

\begin{remark}
Pilih suatu sistem wakil berbeda $X$ untuk $G/H$ dan
tuliskan $\operatorname{Ind}^G_H W = \bigoplus_{x \in X} W$
seperti dalam pembuktian Proposisi~\ref{prop:ind_construction}.
Terdapat pemetaan linear $H$-ekuivarian
\[
\pi : \operatorname{Ind}^G_H W \to W
\]
yang diperoleh dengan mengambil jumlah
\[
\pi\left( (w_x)_{x \in X} \right) :=
\sum_{x \in X} w_x
\]
atas semua $x$ dalam suatu sistem wakil berbeda $X$.
Ternyata pemetaan $\pi$ tidak bergantung pada pilihan $X$
dan dapat dipandang sebagai perumuman teras.

Mendefinisikan pemetaan ini memberikan sifat universal lain:
untuk sebarang representasi $V$ dari $G$ bersama pemetaan $H$-ekuivarian
$\varphi : V \to W$, terdapat pemetaan $G$-ekuivarian tunggal
$\psi : V \to \operatorname{Ind}^G_H W$ sedemikian sehingga
$\pi \circ \psi = \varphi$.
Dalam konteks yang lebih umum, objek-objek yang memenuhi kedua sifat universal
tersebut mungkin tidak berimpit; objek yang memenuhi
konstruksi kedua ini disebut \emph{representasi koterinduksi}.
\end{remark}

\section{Representasi Monomial}

Dalam praktik, banyak representasi diinduksi dari subgrup. Kita telah
melihat banyak contohnya. Secara khusus:

\begin{proposition}
Jika $V$ adalah representasi permutasi dari grup hingga $G$,
maka $V$ merupakan jumlah langsung representasi yang diinduksi dari representasi
trivial suatu subgrup.
\end{proposition}

\begin{proof}
Ingat bahwa $V$ memiliki basis $\{e_x\}_{x \in X}$, dengan $X$ suatu himpunan-$G$.
Karena semua himpunan-$G$ adalah gabungan saling lepas dari himpunan-$G$ transitif,
kita melihat bahwa $V$ merupakan jumlah langsung representasi permutasi
yang bersesuaian dengan himpunan-$G$ transitif. Jadi, tanpa mengurangi keumuman,
kita dapat mengasumsikan $X \cong G/H$ untuk suatu subgrup $H \subseteq G$.

Dengan memilih suatu sistem wakil berbeda $S$ untuk $G/H$,
kita sekarang mengidentifikasi suku-suku langsung $sW$ dalam pembuktian
Proposisi~\ref{prop:ind_construction} dengan ruang yang direntang oleh
unsur-unsur basis $\{ e_{sH} \}$ dalam $V$.
Kita dapat memeriksa bahwa aksi grup $G$ identik dalam kedua kasus.
\end{proof}

\begin{example}
Representasi reguler dari $G$ tepat merupakan
representasi $\operatorname{Ind}^G_1 k$
yang diperoleh dengan induksi dari representasi trivial milik grup
trivial.
\end{example}

\begin{definition}
Representasi $(V,\rho)$ dari grup $G$ disebut \emph{representasi
monomial} atau \emph{representasi permutasi terpuntir} jika
terdapat basis $\{e_1,\ldots,e_n\}$ dari $V$
sedemikian sehingga, untuk semua $1 \le i \le n$ dan $g \in G$,
kita mempunyai
\[ \rho_g(e_i)=\lambda e_j\]
untuk suatu $1 \le j \le n$ dan $\lambda \in k^\times$.
\end{definition}

Representasi monomial adalah representasi yang memiliki basis sedemikian sehingga
matriks-matriks yang bersesuaian mempunyai
tepat satu unsur tak nol pada setiap baris dan kolom.
Representasi permutasi adalah representasi yang unsur tak nolnya
selalu sama dengan satu.

\begin{exercise}
Jika $V$ adalah representasi monomial dari grup hingga $G$,
maka $V$ merupakan jumlah langsung representasi yang diinduksi dari representasi
berdimensi satu milik subgrup-subgrup.
\end{exercise}

\section{Karakter Representasi Terinduksi}

Dalam bagian ini, kita mengasumsikan bahwa medan dasar $k$ adalah $\mathbb{C}$.

\begin{definition}
Untuk suatu fungsi kelas $f : G \to \mathbb{C}$, definisikan
$\operatorname{Res}_H^G f : H \to \mathbb{C}$ sebagai
fungsi kelas yang diperoleh dengan merestriksi domain ke $H$.
\end{definition}

Jelas bahwa karakter dari $\operatorname{Res}_H^G V$
adalah $\operatorname{Res}_H^G \chi_V$.

\begin{definition}
Untuk suatu fungsi kelas $f : H \to \mathbb{C}$, kita membentuk fungsi kelas
baru $\operatorname{Ind}_H^G f : G \to \mathbb{C}$ dengan rumus
\[
(\operatorname{Ind}_H^G f)(g) =
\frac{1}{|H|}\sum_{\substack{a \in G\\a^{-1}ga \in H}} f(a^{-1}ga)
\]
untuk $g \in G$.
\end{definition}

\begin{proposition}
Karakter dari $\operatorname{Ind}_H^G V$
adalah $\operatorname{Ind}_H^G \chi_V$.
\end{proposition}

\begin{proof}
Kita menggunakan notasi matriks blok dari
Proposisi~\ref{prop:ind_construction}
dan memilih $X$ sebagai sistem wakil berbeda untuk $G/H$.
Kita ingin menghitung teras $\rho_g$ untuk suatu $g \in G$.
Ini ekuivalen dengan menjumlahkan teras unsur-unsur
$\rho_g^{x,x}$ untuk semua $x \in X$.
Perhatikan bahwa $\rho_g^{x,x} = 0$, kecuali jika $gx=xh$ untuk suatu
$h \in H$.
Jadi, penjumlahan hanya dilakukan atas $x \in X$ sedemikian sehingga
$x^{-1}gx \in H$. Dalam kasus ini, kita mempunyai
$\rho_g^{x,x}=\phi_x \circ \sigma_h \circ \phi_x^{-1}$,
yang memiliki teras sama dengan $\sigma_h$, dengan $h=x^{-1}gx$.
Dengan demikian, kita memperoleh rumus
\[
(\operatorname{Ind}_H^G f)(g) =
\sum_{\substack{x \in X\\x^{-1}gx \in H}} f(x^{-1}gx) .
\]
Karena $f$ adalah fungsi kelas pada $H$,
kita mempunyai $f(x^{-1}gx)=f(y^{-1}gy)$ apabila $xH=yH$.
Jadi, menjumlahkan atas $G$ alih-alih atas $X$ sama dengan
mengalikan dengan kardinalitas koset, yaitu $|H|$.
\end{proof}

\begin{exercise}
Buktikan bahwa fungsi
$\operatorname{Res}_H^G : C(G) \to C(H)$
adalah homomorfisme gelanggang.
\end{exercise}

\begin{exercise}
Buktikan bahwa fungsi
$\operatorname{Ind}_H^G : C(H) \to C(G)$
adalah homomorfisme grup aditif,
tetapi merupakan homomorfisme gelanggang hanya jika $G=H$.
(Petunjuk: berapakah $(\operatorname{Ind}_H^G 1)(1)$?)
\end{exercise}

\begin{proposition}
Jika $\varphi$ adalah fungsi kelas pada $H$ dan $\psi$ fungsi kelas
pada $G$, maka
$\operatorname{Ind}^G_H ( (\operatorname{Res^G_H} \psi) \cdot\varphi)
= \psi \cdot (\operatorname{Ind}^G_H \varphi)$.
\end{proposition}

\begin{proof}
Karena $\psi(a^{-1}ga)=\psi(g)$ untuk semua $a \in G$, kita mempunyai
\[
\frac{1}{|H|}\sum_{\substack{a \in G\\a^{-1}ga \in H}} \psi(a^{-1}ga)\varphi(a^{-1}ga)
 =
 \psi(g) \left(\frac{1}{|H|}\sum_{\substack{a \in G\\a^{-1}ga \in H}} \varphi(a^{-1}ga)
\right)
\]
untuk semua $g \in G$.
\end{proof}

\begin{corollary}
Bayangan dari $\operatorname{Ind}^G_H$ merupakan ideal dalam $C(G)$.
\end{corollary}

Fakta-fakta di atas dinyatakan untuk gelanggang fungsi kelas
$C(G)$, tetapi dapat diperiksa bahwa fakta-fakta tersebut juga berlaku untuk
gelanggang representasi $R(G)$.

\section{Resiprositas Frobenius}

\begin{theorem}
Misalkan $V$ adalah representasi dari $G$
dan $W$ representasi dari subgrup $H$ dari $G$.
Terdapat isomorfisme kanonik
\[
\operatorname{Hom}^G_k(\operatorname{Ind}_H^G W, V) =
\operatorname{Hom}^H_k(W,\operatorname{Res}_H^G  V).
\]
\end{theorem}

\begin{proof}
Ingat bahwa kita memiliki pemetaan $H$-ekuivarian
$\iota : W \to \operatorname{Ind}_H^G(W)$ menurut definisi
representasi terinduksi.
Jadi, dari suatu unsur
$\psi \in \operatorname{Hom}^G_k(\operatorname{Ind}_H^G W, V)$,
kita memperoleh unsur
$\varphi \in \operatorname{Hom}^H_k(W,\operatorname{Res}_H^G  V)$
dengan komposisi $\varphi = \psi \circ \iota$.
Sifat pemetaan universal memungkinkan kita memperoleh kembali $\psi$ secara tunggal dari $\varphi$,
sehingga kedua himpunan tersebut berkorespondensi secara bijektif.
Kita telah melihat bahwa komposisi pemetaan linear menginduksi transformasi linear
pada ruang-ruang Hom; dengan demikian, kita memperoleh isomorfisme
ruang vektor yang diinginkan.
\end{proof}

Korolari langsungnya adalah hasil terkenal berikut:

\begin{theorem}[Resiprositas Frobenius]
Jika $\varphi$ adalah fungsi kelas pada $H$ dan $\psi$ fungsi kelas
pada $G$, maka
\[
(\operatorname{Ind}_H^G \varphi, \psi)_G
=
(\varphi, \operatorname{Res}_H^G \psi)_H
\]
di mana subskrip digunakan untuk menegaskan bahwa hasil kali dalam pertama berada dalam $C(G)$,
sedangkan hasil kali dalam kedua berada dalam $C(H)$.
\end{theorem}

\begin{proof}
Ketika $\varphi$ dan $\psi$ adalah karakter,
kita mempunyai
\begin{align*}
&(\operatorname{Ind}_H^G \varphi, \psi)_G\\
=& \dim_\mathbb{C}\ \operatorname{Hom}^G_\mathbb{C}(\operatorname{Ind}_H^G W, V)\\
=& \dim_\mathbb{C}\ \operatorname{Hom}^H_\mathbb{C}(W,\operatorname{Res}_H^G  V)\\
=&(\varphi, \operatorname{Res}_H^G \psi)_H .
\end{align*}
Melalui linearitas, hal ini meluas ke semua fungsi kelas karena karakter-karakter
merentang ruang tersebut.
\end{proof}

\section{Dekomposisi Mackey}

Di sini kita mempelajari interaksi antara restriksi dan induksi.

\begin{theorem}[Dekomposisi Mackey]
Misalkan $K,H$ adalah subgrup dari $G$ dan $(W,\sigma)$ suatu representasi dari
$H$.
Maka
\[
\operatorname{Res}^G_K \operatorname{Ind}^G_H W
\cong
\bigoplus_{[x] \in K\backslash G / H}
\operatorname{Ind}_{H_x}^K W_x
\]
di mana $H_x = xHx^{-1} \cap K$, dan $(W_x,\sigma_x)$ adalah representasi
dari $H_x$ dengan ruang vektor dasar $W$
yang didefinisikan oleh $\sigma_x(g)=\sigma(x^{-1}gx)$ untuk semua $g \in H_x$.
\end{theorem}

\begin{proof}
Misalkan $X$ adalah sistem wakil berbeda untuk $G/H$.
Perhatikan bahwa $H_x$ merupakan penstabil dari aksi $K$ pada setiap $x \in X$.
Misalkan $X_1,\ldots,X_r$ adalah subhimpunan-subhimpunan $X$ yang bersesuaian dengan
orbit-orbit $K$ pada $G/H$.
Perhatikan bahwa himpunan $\{X_i\}$ berkorespondensi secara bijektif dengan
koset-koset ganda dalam $K \backslash G / H$.

Kita mengingat deskripsi eksplisit $\tau = \operatorname{Ind}^G_H \rho$
dari Proposisi~\ref{prop:ind_construction}
sebagai jumlah langsung $\bigoplus_{x \in X} xW$,
dengan isomorfisme $\phi_x: W \to xW$ dan $\tau_g$ yang terurai sebagai
\[
\rho_g^{y,x} =
\begin{cases}
\phi_y \circ \rho_{y^{-1}gx} \circ \phi_x^{-1} & \textrm{jika $gxH=yH$,}\\
0 & \textrm{selain itu},
\end{cases}
\]
dalam setiap $\operatorname{Hom}_k(xW,yW)$
untuk $x,y \in X$ dan $g \in G$.

Misalkan $W_i$ adalah jumlah subruang-subruang $xW$ untuk $x \in X_i$.
Kita mempunyai
\[
\operatorname{Ind}^G_H W = \bigoplus_{i=1}^r W_i
\]
di mana setiap $W_i$ mudah dilihat stabil terhadap $K$.
Masih perlu dibuktikan bahwa
$W_i \cong \operatorname{Ind}_{H_x}^K W_x$ untuk suatu pilihan $x \in X_i$.

Sekali lagi menurut Proposisi~\ref{prop:ind_construction},
kita mempunyai $\tau_i = \operatorname{Ind}_{H_x}^K \sigma_x$
dengan $\psi_y : W_x \to yW_x$ dan
\[
\rho_g^{z,y} =
\begin{cases}
\psi_z \circ \rho_{z^{-1}gy} \circ \psi_y^{-1} & \textrm{jika $gyH_x=zH_x$,}\\
0 & \textrm{selain itu},
\end{cases}
\]
dalam setiap $\operatorname{Hom}_k(yW,zW)$
untuk $y,z \in X_i$ dan $g \in K$.

Untuk setiap $y \in X_i$, definisikan $\pi_y : yW \to yW_x$
melalui $\pi_y = \psi_y \circ \phi_y^{-1}$.
Jumlah langsung $\bigoplus_{y \in X_i} \pi_y$
merupakan isomorfisme dari $W_i$ ke $\operatorname{Ind}_{H_x}^K \sigma_x$.
Dapat diperiksa bahwa pemetaan ini $K$-ekuivarian, seperti yang diinginkan.
\end{proof}

Sekarang kita beralih ke penerapan ketika $H=K$.
Seperti dalam pernyataan Teorema Dekomposisi Mackey,
untuk setiap unsur $x \in G$, kita definisikan $H_x = xHx^{-1} \cap H$.
Perhatikan bahwa, meskipun $\operatorname{Res}^H_{H_x} W$
dan $W_x$ sama-sama merupakan representasi dari $H_x$,
keduanya secara umum tidak isomorfik!
Jika $\chi$ adalah karakter dari $W$ dalam $C(H)$, misalkan $\chi_x$ menyatakan
karakter dari $W_x$ dalam $C(H_x)$.

\begin{theorem}[Kriteria Ketaktereduksian Mackey]
Representasi $\operatorname{Ind}^G_H W$ tak tereduksi
jika dan hanya jika $W$ tak tereduksi dan,
untuk setiap $x \in G \setminus H$,
kita mempunyai $( \chi_x, \operatorname{Res}^H_{H_x} \chi )=0$.
\end{theorem}

\begin{proof}
Kita menggunakan Resiprositas Frobenius dan Dekomposisi Mackey berulang kali:
\begin{align*}
&( \operatorname{Ind}^G_H \chi, \operatorname{Ind}^G_H \chi )\\
=&( \chi, \operatorname{Res}^G_H \operatorname{Ind}^G_H \chi )\\
=&( \chi, \sum_{[x] \in H\backslash G / H}
\operatorname{Ind}_{H_x}^H \chi_x )\\
=&\sum_{[x] \in H\backslash G / H} ( \chi, 
\operatorname{Ind}_{H_x}^H \chi_x )\\
=&\sum_{[x] \in H\backslash G / H} ( \operatorname{Res}_{H_x}^H\chi, 
 \chi_x ).
\end{align*}
Jika $x \in H$, maka keduanya sama dan
$\chi=\operatorname{Res}_{H_x}^H\chi = \chi_x$.
Jadi,
\[
( \operatorname{Ind}^G_H \chi, \operatorname{Ind}^G_H \chi )
=(\chi,\chi)+
\sum_{\substack{[x] \in H\backslash G / H\\x \not\in  H}}
( \operatorname{Res}_{H_x}^H\chi,  \chi_x )
\]
Karena $\operatorname{Res}_{H_x}^H\chi$ dan $\chi_x$ adalah karakter sejati
dari $H_x$, hasil kali dalamnya selalu merupakan bilangan bulat tak negatif.
Selain itu, $(\chi,\chi)$ adalah bilangan bulat positif.
Jadi,
$( \operatorname{Ind}^G_H \chi, \operatorname{Ind}^G_H \chi )=1$
jika dan hanya jika $(\chi,\chi)=1$ dan setiap suku langsung
lainnya lenyap.
Teorema pun mengikuti.
\end{proof}

\begin{corollary}
Misalkan $H$ normal dalam $G$.
Representasi $\operatorname{Ind}^G_H \rho$ tak tereduksi
jika dan hanya jika $\rho$ tak tereduksi dan $\rho_x \ne \rho$
untuk semua $x \not\in H$.
\end{corollary}

\begin{proof}
Jika $H$ normal, $H_x=H$ karena $xHx^{-1}=H$ untuk semua $x \in G$.
Secara khusus, $\operatorname{Res}^G_H(\rho)=\rho$.
Dengan mengasumsikan $\chi$ tak tereduksi, syarat
$\left(\chi_x,\operatorname{Res}^G_H(\chi)\right)=0$
ekuivalen dengan $\chi_x \ne \chi$, dan hasilnya mengikuti dari
Kriteria Ketaktereduksian Mackey.
\end{proof}


\section{Contoh}

\subsection{Grup Dihedral}

Misalkan $G$ adalah grup dihedral
\[
D_{2n} = \langle r,s \mid s^2, r^n, (sr)^2 \rangle
\]
dan misalkan $H$ adalah subgrup siklik $\langle r \rangle$.
Misalkan $(W,\sigma)$ adalah representasi berdimensi $1$ dari $H$
yang diberikan oleh $\sigma(r)=(\zeta)$, dengan $\zeta$ suatu akar kesatuan berorde $n$.

Kita akan membentuk $\operatorname{Ind}^G_H W$, yang kita sebut $(V,\rho)$
demi keringkasan.
Kita menggunakan notasi dari Proposisi~\ref{prop:ind_construction}.
Misalkan $X = \{e, s \}$ adalah suatu sistem wakil berbeda untuk $G/H$.
Untuk menentukan $\rho$, kita perlu memahami $\rho_g^{x,y}$
untuk pembangkit $g=r,s$ dan semua pasangan $x,y \in X$.
Setelah menghitung
\[
r(eH)=eH,\ r(sH)=sH,\ s(eH)=sH,\ r(eH)=rH 
\]
kita memperoleh
\[
\rho(r)=
\begin{pmatrix} \sigma(e^{-1}re) & 0 \\ 0 & \sigma(s^{-1}rs) \end{pmatrix}
=\begin{pmatrix} \zeta & 0 \\ 0 & \zeta^{-1} \end{pmatrix}
\]
dan
\[
\rho(s)=
\begin{pmatrix} 0 & \sigma(e^{-1}ss) \\ \sigma(s^{-1}se) & 0 \end{pmatrix}
=\begin{pmatrix} 0 & 1 \\ 1 & 0 \end{pmatrix}.
\]

\subsection{Grup Simetris pada $4$ Huruf}

Misalkan $G$ adalah grup simetris $S_4$ pada $4$ huruf,
dan misalkan $H$ subgrup yang dibangkitkan oleh $(1\ 2)(3\ 4)$
dan $(2\ 3)$.
Perhatikan bahwa $H$ adalah subgrup tak normal berindeks $3$
yang isomorfik dengan $D_8$.
Misalkan $(W,\sigma)$ adalah representasi berdimensi $1$ dari $H$
yang diberikan oleh $\sigma((1\ 2)(3\ 4))=(-1)$ dan $\sigma((2\ 3))=(1)$.
Perhatikan bahwa kernel dari $\sigma$ dibangkitkan oleh $(1\ 4)$ dan $(2\ 3)$.

Seperti sebelumnya, kita membentuk representasi terinduksi $(V,\rho)$ menggunakan
Proposisi~\ref{prop:ind_construction}.
Pertama, ambil $X=\{ e, (1\ 2\ 3), (1\ 3\ 2) \}$ sebagai sistem
wakil berbeda untuk $G/H$.
Kita menghitung
\begin{align*}
(1\ 2)(3\ 4)&\mapsto
\begin{pmatrix}
-1 & 0 & 0\\
0 & 1 & 0\\
0 & 0 & -1\\
\end{pmatrix}\\
(2\ 3)&\mapsto
\begin{pmatrix}
1 & 0 & 0\\
0 & 0 & 1\\
0 & 1 & 0\\
\end{pmatrix}\\
(1\ 2\ 3)&\mapsto
\begin{pmatrix}
0 & 0 & 1\\
1 & 0 & 0\\
0 & 1 & 0\\
\end{pmatrix}
\end{align*}
dalam kasus ini.

\subsection{Representasi yang Tidak Terinduksi}

Di sini kita mengembangkan contoh nonabelian berupa representasi tak tereduksi
yang \emph{bukan} representasi terinduksi.
Pada pembacaan pertama, saya menyarankan untuk sekadar membaca sepintas bagian ini.

Ingat bahwa \emph{aljabar kuaternion Hamilton} $\mathbb{H}$
adalah aljabar pembagian atas $\mathbb{R}$
dengan basis $\{1,i,j,k\}$ dan perkalian yang ditentukan oleh
$ij=k$, $jk=i$, $ki=j$, dan $i^2=j^2=k^2=-1$.
Kita memiliki pembenaman
$\mathbb{H} \hookrightarrow \operatorname{M}_2(\mathbb{C})$
yang diperoleh dengan memperluas pemetaan
\[
1 \mapsto \begin{pmatrix} 1&0\\0&1 \end{pmatrix} \quad
i \mapsto \begin{pmatrix} i&0\\0&-i \end{pmatrix} \quad
j \mapsto \begin{pmatrix} 0&1\\-1&0 \end{pmatrix} \quad
k \mapsto \begin{pmatrix} 0&i\\i&0 \end{pmatrix}
\]
secara linear.
Pembenaman ini merupakan homomorfisme aljabar-$\mathbb{R}$.
(Perhatikan bahwa tidak ada pembenaman $\mathbb{H}$ ke dalam
matriks \emph{real} $2\times 2$.)

Unsur umum $x$ dari $\mathbb{H}$ berbentuk
\[
x = a + bi + cj + dk
\]
dengan $a,b,c,d \in \mathbb{R}$.
Kita mendefinisikan \emph{konjugat} dari $x$ sebagai
\[
\overline{x} = a -bi-cj-dk,
\]
dan \emph{norma} dari $x$ sebagai
\[
\|x\| = \sqrt{\overline{x} x} = \sqrt{a^2+b^2+c^2+d^2},
\]
yang merupakan bilangan real tak negatif.
Kita melihat bahwa $\mathbb{H}$ adalah aljabar pembagian
karena $x^{-1} = \overline{x}/\|x\|$ merupakan rumus eksplisit untuk
invers ketika $x \ne 0$.

Suatu unsur umum $x$ dari $\mathbb{H}$
disebut \emph{kuaternion Lipschitz} jika
$a,b,c,d \in \mathbb{Z}$.
Unsur $x$ disebut \emph{kuaternion Hurwitz}
jika $a,b,c,d$ semuanya bilangan bulat atau semuanya setengah bilangan bulat.
Himpunan $L$ kuaternion Lipschitz dan himpunan $H$ kuaternion Hurwitz
keduanya merupakan subgelanggang dari kuaternion Hamilton.
Grup satuan gelanggang-gelanggang ini dapat ditemukan dengan mengamati bahwa norma
suatu unsur yang dapat dibalik haruslah $1$.

\emph{Grup kuaternion} $Q_8 = \{\pm 1, \pm i, \pm j, \pm k\}$ adalah
grup satuan $U(L)$ dari gelanggang $L$.
Dari pembenaman
$\mathbb{H} \to \operatorname{M}_2(\mathbb{C})$ yang dijelaskan di atas, jelas bahwa
$Q_8$ mempunyai representasi oleh ``matriks permutasi
terpuntir''.
Dengan kata lain, terdapat pilihan basis yang membuat grup bertindak dengan mempermutasikan
unsur-unsur basis dan mengalikannya dengan skalar.

Grup satuan $U(H)$ dari gelanggang $H$ memuat $Q_8$ dan 16 unsur
\[
\frac{ \pm 1 \pm i \pm j \pm k }{2}
\]
dengan semua kombinasi tanda yang mungkin diperbolehkan.
Dapat diperiksa bahwa $\frac{1}{2}(1+i+j+k)$ berorde $6$, sehingga kita memperoleh
representasi dengan bayangan
\[
\left\langle \begin{pmatrix} i&0\\0&-i \end{pmatrix},
\begin{pmatrix} 0&1\\-1&0 \end{pmatrix},
\frac{1}{2} \begin{pmatrix} 1+i&1+i\\-1+i&1-i \end{pmatrix}\right\rangle
.
\]

\begin{exercise}
Buktikan bahwa $U(H)$ tidak memiliki subgrup berindeks $2$.
\end{exercise}

Konsekuensi langsung latihan ini ialah $U(H)$ tidak mungkin
diinduksi dari representasi berdimensi $1$.

Ingat bahwa $A_4$ sering disebut ``grup tetrahedral''
karena isomorfik dengan himpunan matriks ortogonal yang mempertahankan suatu
tetrahedron beraturan.
Oleh karena itu, grup $U(H)$ kadang-kadang disebut ``grup tetrahedral biner''
karena latihan berikut.

\begin{exercise}
Misalkan $Z$ adalah pusat dari $U(H)$.
Buktikan bahwa pusat $Z$ dari $U(H)$ berorde $2$ dan
$U(H)/Z $ isomorfik dengan $A_4$.
\end{exercise}


\begin{exercise}
Buktikan bahwa $U(H)$ isomorfik dengan $Q_8 \rtimes (\mathbb{Z}/3\mathbb{Z})$.
\end{exercise}

\begin{exercise}
Buktikan bahwa $U(H)$ isomorfik dengan grup linear khusus
$\operatorname{SL}_2(\mathbb{F}_3)$ atas medan dengan $3$ unsur.
\end{exercise}

\section{Penerapan}

Misalkan $G$ adalah grup hingga dan $\phi: G \to G$
suatu automorfisme grup. Untuk setiap fungsi kelas $\chi$,
definisikan fungsi kelas $\phi \cdot \chi := \chi \circ \phi^{-1}$.

\begin{exercise}
Buktikan pernyataan-pernyataan berikut.
Pemetaan $\chi \mapsto \phi \cdot \chi$ menginduksi suatu aksi
dari $\operatorname{Aut}(G)$ pada $C(G)$.
Aksi tersebut mempertahankan struktur gelanggang maupun struktur ruang hasil kali dalam pada
$G$, memetakan karakter ke karakter,
dan memetakan karakter tak tereduksi ke karakter tak tereduksi.
Subgrup $\operatorname{Inn}(G)$ dari automorfisme dalam termuat dalam
kernel aksi tersebut.
\end{exercise}

\begin{lemma} \label{lem:normal_restriction}
Misalkan $N$ adalah subgrup normal dari $G$ dan $V$
suatu representasi tak tereduksi dari $G$. Salah satu pernyataan berikut berlaku:
\begin{enumerate}
\item $\operatorname{Res}^G_N V$ isotipik, atau
\item $V \cong \operatorname{Ind}^G_H W$, di mana $W$ adalah representasi tak tereduksi
dari suatu subgrup sejati $H$ dari $G$ sedemikian sehingga $N
\subseteq H$.
\end{enumerate}
\end{lemma}

\begin{proof}
Tinjau dekomposisi isotipik
\[
Res^G_N V \cong U_1 \oplus U_2 \oplus \cdots \oplus U_r
\]
di mana $U_1, \ldots, U_r$ adalah subrepresentasi-subrepresentasi isotipik berbeda dari $N$.
Jika $r=1$, kita berada dalam kasus pertama; jadi kita dapat mengasumsikan $r \ge 2$.
Untuk sebarang $g \in G$ dan indeks $i$, kita melihat bahwa $g(U_i)$ merupakan subruang
$N$-stabil dari $V$; memang, untuk setiap $n\in N$ terdapat $n' \in N$
sedemikian sehingga $ngU_i=gn'U_i=gU_i$.
Misalkan $H$ adalah himpunan semua $g \in G$ sedemikian sehingga $g(U_1)=U_1$.
Jadi, $U$ merupakan subrepresentasi-$H$ dari $V$.
Menurut sifat pemetaan universal, terdapat morfisme $G$-ekuivarian
$\operatorname{Ind}^G_H U_1 \to V$.
Karena $G$ bertindak secara transitif pada semua $U_i$, kita memperoleh isomorfisme.
\end{proof}

\begin{definition}
Grup hingga $G$ disebut \emph{supersolubel} jika terdapat rantai
subgrup
\[
1=G_0 \subseteq G_1 \subseteq \cdots \subseteq G_r = G
\]
di mana setiap $G_i$ normal dalam $G$ dan $G_{i+1}/G_i$ siklik
untuk setiap $i=0,\ldots, r$.
\end{definition}

\begin{exercise}
Buktikan implikasi-implikasi berikut untuk grup hingga $G$:
\[
\textrm{grup-$p$} \implies
\textrm{nilpoten} \implies \textrm{supersolubel} \implies
\textrm{solubel},
\]
dan berikan contoh eksplisit yang menunjukkan bahwa konversnya tidak selalu berlaku.
\end{exercise}

\begin{exercise}
Semua subgrup dan grup hasil bagi dari suatu grup supersolubel juga
supersolubel.
\end{exercise}

\begin{lemma} \label{lem:ss_abel}
Jika $G$ adalah grup supersolubel nonabelian, terdapat subgrup abelian normal
$A$ dari $G$ sedemikian sehingga $A$ tidak termuat dalam pusat
dari $G$.
\end{lemma}

\begin{proof}
Misalkan $Z$ adalah pusat dari $G$. Maka $G/Z$ supersolubel dan
terdapat subgrup siklik normal $C \subseteq G/Z$. Misalkan $A$ adalah
prabayangan $C$ melalui pemetaan hasil bagi $G \to G/Z$.
Karena $C$ normal dalam $G/Z$, kita melihat bahwa $A$ normal dalam $G$.
Karena $A$ merupakan ekstensi pusat dengan hasil bagi siklik, subgrup tersebut harus
abelian.
\end{proof}

\begin{theorem} \label{thm:ss_monomial}
Setiap representasi dari grup supersolubel adalah monomial.
\end{theorem}

\begin{proof}
Kita menggunakan induksi pada orde grup $G$.
Tanpa mengurangi keumuman, kita dapat mengasumsikan bahwa representasi
$(V,\rho)$ tak tereduksi.
Karena grup hasil bagi dari grup supersolubel bersifat supersolubel,
berdasarkan induksi kita dapat mengasumsikan bahwa $V$ merupakan representasi setia.
Jika $G$ abelian, pembuktian selesai.
Jadi, kita dapat mengasumsikan bahwa $G$ nonabelian dan, menurut Lema~\ref{lem:ss_abel},
terdapat subgrup abelian normal $A$ dari $G$ yang tidak termuat dalam pusat.

Kita mengklaim bahwa $\operatorname{Res}^G_A V$ \emph{tidak} isotipik.
Representasi isotipik dari grup abelian terdiri atas matriks-matriks
skalar. Karena matriks-matriks tersebut komutatif dengan setiap matriks lain dalam grup,
dan representasinya setia, ini berarti $A$ berada dalam pusat
$G$, suatu kontradiksi.
Jadi, menurut Lema~\ref{lem:normal_restriction},
kita dapat menyimpulkan bahwa $V$ diinduksi dari suatu subgrup sejati $H$.
Karena $H$ supersolubel, hipotesis induksi menyatakan bahwa
$V$ diinduksi dari representasi monomial.
Jadi, $V$ monomial, seperti yang diinginkan.
\end{proof}

\section{Hasil Kali Semidirek}

Sepanjang bagian ini, $G = A \rtimes H$,
dengan $A$ grup abelian hingga dan $H$ grup hingga.

Ingat bahwa $A^\vee :=
\operatorname{Hom}_{\mathrm{grp}}(A,\mathbb{C}^\times)$
adalah grup abelian yang isomorfik dengan $A$ sendiri (meskipun tidak
secara kanonik).
Terdapat aksi kiri dari $G$ pada $A^\vee$ melalui
\[
(g\chi)(a)=\chi(g^{-1}ag)
\]
dengan $g \in G$, $\chi \in A^\vee$, dan $a \in A$.
Karena $A$ abelian, aksi $G$ tersebut memfaktor melalui proyeksi
ke subgrup $H$.

Misalkan $\chi_1,\ldots, \chi_r$ adalah sistem wakil berbeda
untuk orbit-orbit $H$ dari $A^\vee$. Definisikan
$H_i=\operatorname{Stab}_H(\chi_i)$.
Secara khusus,
\[
A^\vee \cong H/H_1 \sqcup \cdots \sqcup H/H_r
\]
sebagai himpunan-$H$.

Untuk suatu representasi tak tereduksi $\rho$ dari $H_i$,
kita akan mendefinisikan representasi $\theta_{i,\rho}$ sebagai berikut.
Misalkan $G_i= A \rtimes H_i$.
Perluas $\rho$ menjadi representasi $\widetilde{\rho}$ dari $G_i$
dengan mengomposisikannya dengan pemetaan $G_i \to H_i$.
Perluas $\chi_i$ menjadi representasi $\widetilde{\chi_i}$ dari $G_i$
dengan mendefinisikan $\widetilde{\chi_i}(ah)=\chi_i(a)$
untuk semua $a \in A$ dan $h \in H_i$; ini bermakna karena
$H_i$ bertindak secara trivial pada $\chi_i$.
Terakhir, kita definisikan
\[
\theta_{i,\rho} := \operatorname{Ind}_{G_i}^G
\widetilde{\chi_i} \otimes \widetilde{\rho},
\]
yang merupakan representasi dari $G$.

\begin{theorem}
Setiap representasi tak tereduksi dari $G$ isomorfik dengan suatu
$\theta_{i,\rho}$,
setiap $\theta_{i,\rho}$ tak tereduksi, dan
$\theta_{i,\rho} \cong \theta_{i', \rho'}$ jika dan hanya jika
$i=i'$ dan $\rho \cong \rho'$.
\end{theorem}

\begin{proof}
Pertama, kita tunjukkan bahwa $\theta_{i,\rho}$ tak tereduksi.
Misalkan $\eta = \widetilde{\chi_i} \otimes \widetilde{\rho}$,
dengan $\theta_{i,\rho} = \operatorname{Ind}_{G_i}^G \eta$.
Misalkan $x \in G \setminus G_i$ dan $K_x = G_i \cap xG_ix^{-1}$.
Perhatikan bahwa $\operatorname{Res}^{K_x}_A \eta = n\chi_i$,
dengan $n$ derajat dari $\eta$.
Seperti dalam pernyataan Dekomposisi Mackey,
kita mendefinisikan representasi dari $K_x$ melalui
$\eta_x (g) = \eta(x^{-1}gx)$ untuk semua $g \in K_x$.
Sekarang, $\operatorname{Res}^{K_x}_A \eta = n\chi'$
dengan $\chi'(a) = \chi_i(x^{-1}ax)$ untuk semua $a \in A$.
Karena $\chi_i(a) \ne \chi_i(x^{-1}ax)$ untuk suatu $a \in A$,
kita menyimpulkan bahwa $\chi_i \ne \chi'$.
Karena $\chi_i$ dan $\chi'$ tak tereduksi, kita menyimpulkan bahwa
$(\chi_i, \chi')=0$. Jadi, $(\eta,\eta_x)=0$, dan kita
menyimpulkan bahwa $\theta_{i,\rho}$ tak tereduksi menurut Kriteria
Ketaktereduksian Mackey.

Sekarang kita tunjukkan bahwa kelas isomorfisme dari $\theta_{i,\rho}$
ditentukan oleh $i$ dan $\rho$.
Restriksi $\operatorname{Res}^G_A \theta_{i,\rho}$
merupakan jumlah langsung karakter-karakter $\nu_1,\ldots,\nu_s$ dalam $A^\vee$.
Dari Teorema Dekomposisi Mackey, kita mengetahui bahwa setiap
$\nu_j$ berbentuk $g \cdot \chi_i$ untuk suatu $g \in G$.
Karena $A$ abelian, setiap $\nu_j$ berbentuk $h \cdot \chi_i$
untuk suatu $h \in H$. Secara khusus, semua $\nu_j$ berasal dari orbit-$H$
yang sama dalam $A^\vee$, sehingga indeks $i$ ditentukan.
Misalkan $W$ adalah komponen isotipik dari $\operatorname{Res}^G_A V$
yang terkait dengan representasi $\chi_i$.
Grup $G_i$ membiarkan $W$ invarian, dan kita melihat bahwa
$W$ isomorfik dengan representasi
$\widetilde{\chi_i} \otimes \widetilde{\rho}$,
yang menentukan $\rho$.

Terakhir, kita tunjukkan bahwa setiap representasi tak tereduksi dari $G$ memiliki
bentuk yang diinginkan. Misalkan $(V, \sigma)$ adalah representasi tak tereduksi dari $G$.
Misalkan
\[
\operatorname{Res}^G_A V = \bigoplus_{\chi \in A^\vee} W_\chi
\]
adalah dekomposisi isotipiknya. Pilih suatu karakter $\chi_i \in A^\vee$
dengan $W_{\chi_i} \ne 0$, dan misalkan $H_i = \operatorname{Stab}_H(\chi_i)$.
Perhatikan bahwa $W_{\chi_i}$ merupakan subrepresentasi dari $V$
untuk grup $H_i$; misalkan $\tau$ adalah representasi-$H_i$ ini.
Kita melihat bahwa aksi $G_i$ pada $W_{\chi_i}$ merupakan representasi
$\widetilde{\chi_i} \otimes \widetilde{\tau}$.
Dengan menggunakan Resiprositas Frobenius, kita memperoleh rangkaian pertidaksamaan
\[
0 < \left( \widetilde{\chi_i} \otimes \widetilde{\tau},
\operatorname{Res}^G_{G_i} \sigma \right)_{G_i}
= \left( \operatorname{Ind}^G_{G_i} \widetilde{\chi_i} \otimes \widetilde{\tau},
\sigma \right)_{G} \le 1
\]
di mana pertidaksamaan pertama mengikuti karena
$W_{\chi_i}$ merupakan suku langsung dari $V$, sedangkan pertidaksamaan kedua
mengikuti karena $\sigma$ tak tereduksi.
Kita menyimpulkan bahwa $\sigma \cong \theta_{i,\tau}$, seperti yang diinginkan.
\end{proof}

\section{Teorema Artin dan Brauer}

Misalkan $\mathcal{F}$ adalah keluarga subgrup dari grup hingga $G$.
Definisikan pemetaan
\[
\Sigma_{\mathcal{F}} \operatorname{Ind} :
\bigoplus_{H \in \mathcal{F}} R(H) \to R(G)
\]
melalui
\[
(\chi_H)_{H \in \mathcal{F}} \mapsto \sum_{H \in \mathcal{F}}
\operatorname{Ind}_H^G \chi_H .
\]

\begin{theorem}[Teorema Artin]
Jika $\mathcal{F}$ adalah himpunan subgrup siklik dari $G$,
maka $\Sigma_{\mathcal{F}} \operatorname{Ind}$ memiliki kokernel hingga.
\end{theorem}

\begin{proof}
Pertama, misalkan $V$ adalah ruang hasil kali dalam kompleks dan tinjau pemetaan-pemetaan
\[
\Sigma : \bigoplus_{i=1}^n V \to V \textrm{ dan }
\Delta : V \to \bigoplus_{i=1}^n V
\]
yang didefinisikan melalui
\[
\Sigma\left(v_1,\ldots,v_n\right) = \sum_{i=1}^n v_i \textrm{ dan }
\Delta(v) = (v,\ldots,v).
\]
Keduanya merupakan operator linear adjoin, dalam arti bahwa
\[
\left( \Sigma(w), v \right) = \left( w, \Delta(v) \right)
\]
untuk semua $w \in \bigoplus_{i=1}^n V$ dan $v \in V$.

Sekarang, fungsi
\[
\left( \bigoplus_{H \in \mathcal{F}} \operatorname{Res}^G_H \right)
\circ \Delta: C(G) \to \bigoplus_{H \in \mathcal{F}} C(H)
\]
bersifat injektif karena fungsi kelas ditentukan oleh nilai-nilainya pada
subgrup-subgrup siklik.
Jika suatu transformasi linear injektif, operator adjoinnya
\[
\Sigma_{\mathcal{F}} \operatorname{Ind} :
\bigoplus_{H \in \mathcal{F}} C(H) \to C(G)
\]
bersifat surjektif.

Sekarang, pemetaan linear
\[
\Sigma_{\mathcal{F}} \operatorname{Ind} :
\mathbb{Q} \otimes_{\mathbb{Z}}
\left(\bigoplus_{H \in \mathcal{F}} R(H)\right)
\to \mathbb{Q} \otimes_{\mathbb{Z}} R(G)
\]
harus surjektif, karena penensoran dengan $\mathbb{C}$ menghasilkan
pemetaan linear surjektif.
Oleh karena itu, restriksinya ke $\bigoplus_{H \in \mathcal{F}} R(H)$
memiliki kokernel hingga, karena bayangannya merupakan subgrup dari $R(G)$
dengan rank maksimal.
\end{proof}

\begin{definition}
Suatu grup $H$ disebut \emph{$p$-elementer} untuk bilangan prima $p$ jika
$H \cong P \times C$, dengan $P$ grup-$p$ dan $C$ grup siklik
yang ordenya relatif prima terhadap $p$.
Kita menyebut grup $H$ \emph{elementer} jika $p$-elementer untuk suatu
bilangan prima $p$.
\end{definition}

\begin{theorem}[Teorema Brauer]
Jika $\mathcal{F}$ adalah himpunan subgrup elementer dari $G$,
maka $\Sigma_{\mathcal{F}} \operatorname{Ind}$ surjektif.
\end{theorem}

\begin{proof}
Lihat \cite[\S{10}]{Serre}.
\end{proof}

\begin{corollary}
Setiap karakter merupakan kombinasi linear bilangan bulat dari karakter-karakter monomial.
\end{corollary}

\begin{proof}
Teorema Brauer menyatakan bahwa setiap karakter merupakan kombinasi linear bilangan bulat
dari karakter-karakter yang diinduksi dari subgrup elementer.
Grup elementer bersifat nilpoten, sehingga supersolubel.
Menurut Teorema~\ref{thm:ss_monomial}, semua representasinya
monomial.
Representasi yang diinduksi dari representasi monomial juga monomial.
\end{proof}

\bibliographystyle{alpha}
\bibliography{rep_theory}

\end{document} 


